%!TEX root =  tfg.tex
\chapter{Arranque}

\begin{quotation}[Novelist]{Ernest Hemingway (1899--1961)}
The good parts of a book may be only something a writer is lucky enough to overhear or it may be the wreck of his whole damn life -- and one is as good as the other.
\end{quotation}

\begin{abstract}
En este capítulo se formaliza el alcance del proyecto y se 
establecen las bases técnicas para su desarrollo. El objetivo es 
traducir la visión conceptual de Via Inferni en una lista estructurada de 
requisitos funcionales, aplicando una primera iteración de la metodología Feature 
Driven Development (FDD). Asimismo, se define la arquitectura de software que soportará
dichas funcionalidades, detallando los entornos de producción, las herramientas de 
diseño de datos (preproducción) y la estrategia de pruebas necesaria para asegurar 
la calidad del software.
\end{abstract}

\section{Lista de características}

Siguiendo la metodología \textit{Feature Driven Development} (FDD), 
se ha descompuesto la visión global del videojuego en conjuntos de características 
agrupados por dominio funcional. Esta lista representa el alcance del 
Producto Mínimo Viable (MVP) y prioriza las mecánicas 
centrales que definen la experiencia \textit{roguelike} en la ambientación de 
la Divina Comedia.

\subsection{Dominio: Jugador y Mecánicas Core}
Este conjunto abarca las funcionalidades críticas relacionadas con el control, 
el combate y el ciclo de vida del avatar.
\begin{itemize}
    \item \textbf{Gestión de Personajes:} Implementación del sistema dual.
    \begin{itemize}
        \item La dualidad de los personaes por definir todavia.
    \end{itemize}
    \item \textbf{Sistema de Combate:}
    \begin{itemize}
        \item Sistema de combate por definir todavia.
    \end{itemize}
    \item \textbf{Ciclo de Vida (Run):}
    \begin{itemize}
        \item Gestionar la salud actual y el estado de muerte (Game Over).
        \item Reiniciar el progreso del mundo al morir (Muerte permanente).
    \end{itemize}
\end{itemize}

\subsection{Dominio: Generación Procedural y Mundo}
Funcionalidades encargadas de la construcción dinámica de los nueve círculos del infierno.
\begin{itemize}
    \item \textbf{Algoritmo de Generación:}
    \begin{itemize}
        \item Generar estructura de cuadrícula de A a B salas aleatorias (cambiar valores).
        \item Seleccionar \textit{layouts} prediseñados de salas (combate, tesoro, tienda, 
        descanso).
        \item Aumentar la complejidad del mapa según la profundidad del círculo.
    \end{itemize}
    \item \textbf{Navegación:}
    \begin{itemize}
        \item Gestionar el estado de las puertas (bloqueo por enemigos vivos).
        \item Transicionar entre niveles (círculos) al derrotar al jefe.
    \end{itemize}
\end{itemize}

\subsection{Dominio: Entidades y Comportamiento}
Definición de los agentes hostiles que pueblan el juego.
\begin{itemize}
    \item \textbf{Arquetipos de Enemigos (10-15 variantes):}
    \begin{itemize}
        \item \textit{Chaser:} Navegación directa hacia el jugador (Pathfinding).
        \item \textit{Ranged:} Cálculo de trayectoria y disparo de proyectiles.
        \item \textit{Tank:} Comportamiento defensivo y alta resistencia.
        \item \textit{Flying:} Movimiento ignorando obstáculos del terreno.
    \end{itemize}
    \item \textbf{Jefes Finales:}
    \begin{itemize}
        \item \textbf{Cancerbero:} Patrones de ataques por definir todavia.
        \item \textbf{Satán:} Patrones de ataques por definir todavia.
    \end{itemize}
\end{itemize}

\subsection{Dominio: Economía y Progresión}
Sistemas que otorgan profundidad estratégica a la partida.
\begin{itemize}
    \item \textbf{Sistema de Ítems:}
    \begin{itemize}
        \item Aplicar modificadores de estadísticas (ítems pasivos).
        \item Ejecutar efectos de un solo uso (consumibles).
    \end{itemize}
    \item \textbf{Economía:}
    \begin{itemize}
        \item Generar y recolectar monedas (\textit{drops}) de enemigos.
        \item Transaccionar compra de mejoras en la sala de tienda.
    \end{itemize}
\end{itemize}

\subsection{Dominio: Interfaz y Feedback (UI/UX)}
\begin{itemize}
    \item \textbf{HUD (Heads-Up Display):} Visualizar vida, personaje activo, 
    minimapa y \textit{cooldowns}.
    \item \textbf{Flujo de Pantallas:} Navegación entre Menú Principal, 
    Pausa y Pantalla de Resultados.
\end{itemize}

\section{Diseño arquitectónico}

\subsection{Visión general del sistema}

El videojuego desarrollado en el presente Trabajo de Fin de Grado consiste en un roguelike con generación procedural de niveles inspirado en la estructura del Infierno de Dante. El juego se organiza en nueve círculos temáticos, cada uno con identidad visual propia, enemigos específicos, conjunto de salas diferenciado y un jefe final que permite el acceso al siguiente círculo.

Desde el punto de vista arquitectónico, el sistema se ha diseñado siguiendo una estructura modular orientada a la separación de responsabilidades y al bajo acoplamiento entre componentes. La arquitectura se basa en un conjunto de gestores especializados (manager-based architecture), donde cada subsistema tiene una responsabilidad claramente delimitada.

Adicionalmente, el sistema contempla la existencia de dos personajes jugables, Dante y Virgilio, cuya lógica de control y estado debe poder alternarse dinámicamente durante la ejecución. Aunque la mecánica de diferenciación entre ambos personajes no se encuentra completamente definida en esta fase, la arquitectura global se ha diseñado para permitir su extensión sin modificaciones estructurales significativas.

\subsection{Arquitectura global del sistema}

El sistema se organiza en los siguientes subsistemas principales:

\begin{itemize}
    \item \textbf{GameManager}: Controlador global del estado del juego.
    \item \textbf{CircleManager}: Gestión del círculo actual.
    \item \textbf{LevelGenerator}: Generación estructural del nivel.
    \item \textbf{RoomManager}: Control de activación de salas.
    \item \textbf{RoomGenerator}: Generación de contenido interno.
    \item \textbf{DifficultyManager}: Ajuste dinámico de dificultad.
    \item \textbf{CharacterManager}: Gestión de personajes jugables.
\end{itemize}

El siguiente diagrama ilustra la estructura general y las relaciones entre subsistemas:

\begin{center}
\begin{tikzpicture}[
    node distance=1.8cm,
    box/.style={rectangle, draw, rounded corners, align=center, minimum width=3cm, minimum height=0.9cm},
    arrow/.style={-Latex}
]

\node[box] (gm) {GameManager};

\node[box, below=of gm] (cm) {CircleManager};
\node[box, left=of cm] (chm) {CharacterManager};
\node[box, right=of cm] (dm) {DifficultyManager};

\node[box, below=of cm] (lg) {LevelGenerator};
\node[box, below=of lg] (rm) {RoomManager};
\node[box, below=of rm] (rg) {RoomGenerator};

\draw[arrow] (gm) -- (cm);
\draw[arrow] (gm) -- (chm);
\draw[arrow] (gm) -- (dm);

\draw[arrow] (cm) -- (lg);
\draw[arrow] (lg) -- (rm);
\draw[arrow] (rm) -- (rg);

\draw[arrow] (dm) -- (rg);

\end{tikzpicture}
\end{center}

\subsection{Gestión de personajes jugables}

El sistema contempla dos personajes jugables: Dante y Virgilio. Para garantizar extensibilidad, se ha diseñado un \textit{CharacterManager} responsable de:

\begin{itemize}
    \item Mantener el estado del personaje activo.
    \item Gestionar el posible cambio dinámico entre personajes.
    \item Centralizar atributos diferenciadores (estadísticas, habilidades, modificadores).
\end{itemize}

Ambos personajes comparten una interfaz común de comportamiento, lo que permite desacoplar la lógica del jugador del resto del sistema. De esta forma, la introducción de mecánicas específicas para cada personaje no afecta a los módulos de generación de niveles o enemigos.

\subsection{Modelo estructural del nivel}

La generación del nivel se basa en una arquitectura híbrida compuesta por:

\begin{itemize}
    \item Un modelo espacial basado en una cuadrícula bidimensional (grid) que controla la ocupación geométrica.
    \item Un modelo lógico basado en grafos donde cada nodo representa una habitación y cada arista una conexión navegable.
\end{itemize}

El grid se emplea exclusivamente para validar la coherencia espacial durante la generación. Cada habitación define un conjunto de celdas ocupadas (footprint), lo que permite soportar geometrías complejas como salas rectangulares o en forma de L.

Sin embargo, desde el punto de vista lógico, cada habitación constituye un único nodo en el grafo, independientemente de su tamaño. Esto permite desacoplar la estructura espacial de la conectividad del nivel.

\subsection{Proceso de generación procedural}

La generación se divide en dos fases claramente diferenciadas.

\subsubsection{Generación estructural}

En esta fase, el \textit{LevelGenerator} construye el grafo completo del círculo actual. El proceso sigue una estrategia de expansión iterativa:

\begin{enumerate}
    \item Se crea la sala inicial.
    \item Se selecciona una puerta libre de una sala existente.
    \item Se elige una plantilla de habitación compatible.
    \item Se verifica su encaje espacial en el grid.
    \item Si el encaje es válido, se actualizan grid y grafo.
\end{enumerate}

Una vez finalizada la generación, se determina la sala del jefe como el nodo más alejado de la sala inicial, utilizando un recorrido en anchura (BFS).

\subsubsection{Generación de contenido bajo demanda}

El contenido interno de las salas no se genera durante la fase estructural. En su lugar, el \textit{RoomGenerator} instancia enemigos, objetos y eventos únicamente cuando el jugador accede a una sala, así como en sus adyacentes inmediatas.

Este enfoque reduce el consumo de memoria y permite que el \textit{DifficultyManager} ajuste dinámicamente la complejidad del contenido en función del estado actual del jugador.

\subsection{Configuración basada en datos}

Cada uno de los nueve círculos se define mediante una estructura de configuración independiente que incluye:

\begin{itemize}
    \item Conjunto de enemigos disponibles.
    \item Conjunto de plantillas de habitaciones.
    \item Conjunto de objetos.
    \item Prefab del jefe.
    \item Parámetros de dificultad.
\end{itemize}

Esta aproximación sigue un paradigma \textit{data-driven}, separando configuración y comportamiento. Esto permite modificar el balance del juego o añadir nuevos círculos sin alterar la lógica central del sistema.

\subsection{Justificación arquitectónica}

Las decisiones adoptadas persiguen los siguientes objetivos:

\begin{itemize}
    \item Separación clara de responsabilidades.
    \item Bajo acoplamiento entre módulos.
    \item Escalabilidad para incorporar nuevas mecánicas.
    \item Soporte para geometrías complejas de habitaciones.
    \item Flexibilidad para integrar múltiples personajes jugables.
\end{itemize}

La combinación de un modelo espacial basado en grid y un modelo lógico basado en grafos permite mantener coherencia geométrica sin sacrificar expresividad estructural. Asimismo, la modularización del sistema facilita su mantenimiento y evolución futura.